\renewcommand{\thechapter}{\Roman{chapter}}
\chapter[]{CONCLUSION}
\renewcommand{\thechapter}{\arabic{chapter}}
\begingroup
\fontsize{12pt}{\baselineskip}\selectfont

\section{Summary of Findings}

\subsection{Effectiveness of the Hierarchical Dual-Stream Tokenizer}
\begin{indentedblock}{4em}{0pt}
	This study successfully demonstrated that translating raw Boundary Representation (B-Rep) data into a deterministic, linearized topological graph is a viable alternative to lossy visual conversions. The Hierarchical Dual-Stream Tokenizer effectively bridged the modality gap between geometric mathematics and semantic language. By utilizing Semantic Macro-Compression, the model reduced sequence complexity while maintaining structural integrity, as evidenced by the rapid convergence of Cross-Entropy loss to approximately 0.5 within 20 epochs. Furthermore, the integration of Rotary Positional Embeddings (RoPE) allowed the model to achieve rotational and translational invariance, enabling it to grasp physical assembly logic directly within the attention mechanism.
\end{indentedblock}

\subsection{Performance Across Hierarchical Annotation Levels}
\begin{indentedblock}{4em}{0pt}
	The proposed architecture, leveraging a Qwen 3.5 backbone, met or exceeded all primary research objectives regarding annotation fidelity. At Level 1, the model achieved a peak BERTScore of 0.95, proving a near-perfect ability to identify functional identities like fasteners and gears. Level 2 structural descriptions reached a ROUGE-L of 0.88, indicating successful reconstruction of complex syntactic flows describing topological features. Most significantly, the model attained a BLEU-4 score of 0.85 for Level 3 parametric specifications[cite: 77, 961]. This demonstrates that the framework can generate precise engineering constraints and dimensions without requiring a construction history, a critical requirement for processing industrial-standard STEP files.
\end{indentedblock}

\section{Significance and Contributions}
\begin{indentedblock}{4em}{0pt}
	Applying the LIMA principle proved that data quality and density are superior to sheer volume for specialized CAD domain adaptation as the small dataset is able to accurately annotate the files reserved for testing. The approach eliminated spurious noise and "shortcut learning" common in models trained on unverified web-scraped data. The resulting framework revolutionizes industrial CAD management by enabling automated, high-fidelity indexing and retrieval, significantly reducing the manual overhead previously required for design documentation.
\end{indentedblock}

\section{Limitations and Future Research}
\begin{indentedblock}{4em}{0pt}
	Despite high quantitative performance and a significant agreement in human verification quantified using Cohen's Kappa, certain limitations remain. For example, the model is currently constrained to the English language and focuses on common reusable assets due to the scarcity of complex, publicly available industrial assemblies. Furthermore, the current model is limited to only generating technical, descriptive annotations in English. Future studies could expand the vocabulary to include multi-lingual annotations and explore zero-shot adaptation for even more specialized engineering domains. Additionally, further research into reducing the computational cost of Graph Isomorphism matching during the macro-compression stage would improve pipeline scalability for large-scale enterprise repositories.
\end{indentedblock}

\endgroup